\clearpage
\section{相关真题：曲线积分与曲面积分}
\src{§5.4，pp.296—297}
本节为筛选摘录。原书例 5.4.1—5.4.14 为纯 Riemann 积分、函数与积分不等式等（包括例 5.4.11 的纯二重积分），均不收入；此处仅收例 5.4.15、5.4.16。

\entry{例 5.4.15}{第九届全国决赛题}{pp.296—297}
求第二型曲线积分
\[
I=\int_\Gamma e^{\sin x}(\cos x\cos y\dd x-\sin y\dd y)+\cos z\dd z,
\]
其中 $\Gamma:x=\pi\sin(t/2)$、$y=t-\sin t$、$z=\sin2t$，$0\le t\le\pi$。
\sol 记 $P=e^{\sin x}\cos x\cos y$、$Q=-e^{\sin x}\sin y$、$R=\cos z$，则 $P,Q,R$ 在 $\mathbb R^3$ 上连续可微，且
\[
\rot(P,Q,R)=(R_y-Q_z,P_z-R_x,Q_x-P_y)=(0,0,0),
\]
故积分仅与端点有关。$\Gamma$ 的端点为 $(0,0,0)$、$(\pi,\pi,0)$，因此
\begin{align*}
I&=\int_{(0,0,0)}^{(\pi,\pi,0)}P\dd x+Q\dd y+R\dd z\\
&=\int_{(0,0,0)}^{(\pi,0,0)}P\dd x+Q\dd y+R\dd z
+\int_{(\pi,0,0)}^{(\pi,\pi,0)}P\dd x+Q\dd y+R\dd z\\
&=\int_0^\pi e^{\sin x}\cos x\dd x+\int_0^\pi(-\sin y)\dd y=0-2=-2.
\end{align*}
\ednote{原解 $P$ 漏 $\cos x$、$Q$ 漏负号及折线路径首项符号已依题面和方向校正；不改变原解法及末答。}

\entry{例 5.4.16}{第八届全国决赛题}{p.297}
求
\[
I=\iint_S\frac{ax\dd y\dd z+(z+a)^2\dd x\dd y}{\sqrt{x^2+y^2+z^2}},
\]
其中 $a>0$ 为常数，$S:z=-\sqrt{a^2-x^2-y^2}$，取向上侧。
\sol 在 $S$ 上 $x^2+y^2+z^2=a^2$，所以
\[
I=\frac1a\iint_S[ax\dd y\dd z+(z+a)^2\dd x\dd y].
\]
补底 $D:\{x^2+y^2\le a^2,z=0\}$，取向上侧，则 $S^{-1}\cup D$ 为闭下半球体 $B$ 的表面，取向外侧。由 Gauss 公式，
\begin{align*}
&\oiint_{S^{-1}\cup D}[ax\dd y\dd z+(z+a)^2\dd x\dd y]\\
&\quad=\iiint_B(3a+2z)\dd V\\
&\quad=3a\cdot\frac{2\pi a^3}{3}
+2\int_0^{2\pi}\int_{\pi/2}^\pi\int_0^a(r\cos\theta)r^2\sin\theta\dd r\dd\theta\dd\varphi\\
&\quad=2\pi a^4-\frac{\pi a^4}{2}=\frac{3\pi a^4}{2}.
\end{align*}
所以
\begin{align*}
\iint_S[ax\dd y\dd z+(z+a)^2\dd x\dd y]
&=-\frac{3\pi a^4}{2}+\iint_D[ax\dd y\dd z+(z+a)^2\dd x\dd y]\\
&=-\frac{3\pi a^4}{2}+\iint_Da^2\dd x\dd y\\
&=-\frac{3\pi a^4}{2}+\pi a^4=-\frac{\pi a^4}{2}.
\end{align*}
故 $I=-\pi a^3/2$。
\ednote{原书补面 $D$ 仅写 $x^2+y^2=a^2$；此处依 Gauss 公式的应用，明确为 $z=0$ 上的闭圆盘。}
